\section{Requirements and Methodology}

The fictional design uses four sensors with half-wavelength spacing and a nominal operating
frequency of \SI{2.45}{\giga\hertz}. The analysis workflow consists of geometry definition,
analytical modeling, synthetic-data generation, covariance estimation, direction-of-arrival
processing, and reporting.

\subsection{Representative Requirements}

\begin{problem}{Example Engineering Requirements}{prob:requirements}
Design a four-channel array-processing demonstrator that accepts complex sensor samples,
forms a sample covariance matrix, estimates a source direction, and produces a concise
validation report. The demonstrator should document assumptions, numerical settings, and
known limitations.
\end{problem}

\begin{definition}{Validation Evidence}{def:validation}
Validation evidence is a reproducible artifact---for example, an equation check, source-code
test, simulation plot, measurement table, or convergence study---that supports a technical
claim made in the report.
\end{definition}

\subsection{Method Summary}

\begin{enumerate}[label=\arabic*.]
    \item Define the physical and numerical parameters.
    \item Generate or acquire complex sensor snapshots.
    \item Estimate the sample covariance matrix.
    \item Apply a subspace or beamforming estimator.
    \item Compare the recovered result with the known test condition.
    \item Record limitations and recommended follow-up verification.
\end{enumerate}
